% Part: incompleteness
% Chapter: incompleteness-provability
% Section: lob-thm

\documentclass[../../../include/open-logic-section]{subfiles}

\begin{document}

\olfileid{inc}{inp}{tar}

\olsection{真理的不可定义性}

\emph{可定义性}的概念依赖于为算术语言给出的形式语义。
我们已经描述过算术语言中的一组公式和语句。
“预期解释”即是将这样的语句读作关于自然数的断言，而该断言可真可假。
令 $\Struct{N}$ 为以 $\Nat$ 为论域、并对算术语言中的符号作标准解释的结构。
那么 $\Sat{N}{!A}$ 就意味着“$!A$ 在标准解释下为真”。

\begin{defn}
自然数上的一个关系$R(x_1,\dots,x_k)$在 $\Struct{N}$ 中是\emph{可定义的}，
当且仅当在算术语言中存在一个公式$!A(x_1,\dots,x_k)$，
使得对所有 $n_1,\dots,n_k$，
$R(n_1,\dots,n_k)$当且仅当$\Sat{N}{!A(\num n_1,\dots,\num
  n_k)}$。
\end{defn}

换句话说，一个关系在 $\Struct{N}$ 中可定义，
当且仅当它在理论 $\Th{TA}$ 中是可表示的，其中$\Th{TA} =
\Setabs{!A}{\Sat{N}{!A}}$是算术真语句的集合。（如果你不能立刻明白这一点，应当回顾相关定义并说服自己事实正是如此。）

\begin{lem}
每个可计算关系都在~$\Struct{N}$ 中可定义。
\end{lem}

\begin{proof}
容易验证，在 $\Th{Q}$ 中表示某关系的公式，也在 $\Struct{N}$ 中定义同一个关系。 
\end{proof}

现在自然要问，反之亦然吗？即，每个在~$\Struct{N}$ 中可定义的关系都是可计算的吗？答案是否定的。例如：

\begin{lem}
停机关系在 $\Struct{N}$ 中可定义。
\end{lem}

\begin{proof}
回忆 克林 范式定理，它断言每个部分可计算函数~$f$ 都有一个指标~$e$，使得 $f(x) =
U(\umin{s}{T(e,x,s)})$ 对所有 $x \in \Nat$ 成立，其中 $U$ 和 $T$ 是原始递归的，因而是全函数。因此，$f(x)$ 有定义（即计算停机）当且仅当存在某个~$s$ 使得 $T(e,x,s)$ 成立。

现在令 $H$ 为停机关系，即
\[
H = \Setabs{\tuple{e,x}}{\lexists[s][T(e, x, s)]}.
\]
设 $!D_T$ 在 $\Struct{N}$ 中定义 $T$。则
\[
H = \Setabs{\tuple{e,x}}{\Sat{N}{\lexists[s][!D_T(\num e, \num x, s)]}},
\]
因此 $\lexists[s][!D_T(z, x, s)]$ 就在 $\Struct{N}$ 中定义了~$H$。 
\end{proof}

\begin{prob}
证明 $Q(n) \defiff n \in \Setabs{\Gn{!A}}{\Th{Q} \Proves !A}$ 在算术中是可定义的。
\end{prob}

那么 $\Th{TA}$ 本身呢？它在算术中是可定义的吗？也就是说，集合 $\Setabs{\Gn{!A}}{\Sat{N}{!A}}$ 在算术中是可定义的吗？塔斯基定理对此给出了否定的回答。

\begin{thm}
\ollabel{thm:tarski}
算术真语句的集合在算术中是不可定义的。
\end{thm}

\begin{proof} 
假设 $!D(x)$ 定义了该集合，即 $\Sat{N}{!A}$ 当且仅当 $\Sat{N}{!D(\gn{!A})}$。根据不动点引理，存在一个公式 $!A$ 使得 $\Th{Q} \Proves !A \liff \lnot !D(\gn{!A})$，从而 $\Sat{N}{!A \liff \lnot !D(\gn{!A})}$。但这样一来，$\Sat{N}{!A}$ 当且仅当 $\Sat{N}{\lnot !D(\gn{!A})}$，这与 $!D(y)$ 本应定义算术真语句集合的事实相矛盾。  
\end{proof}

塔斯基 将这一分析应用到关于真理的更一般的哲学概念上。对于任何语言 $L$，塔斯基 认为，一个关于 $L$ 的充分的真理概念必须对每个语句 $X$ 满足
\begin{quote}
``$X$'' 是真的当且仅当 $X$。
\end{quote}
塔斯基 常被引用的以英语为例的句子是：
\begin{quote}
``雪是白的''是真的当且仅当雪是白的。
\end{quote}
然而，对于任何强到足以表示对角函数的语言，以及该语言中的任何谓词 $T(x)$，我们都能构造出一个满足``$X$ 当且仅当并非 $T(\text{`$X$'})$''的语句 $X$。考虑到我们不会希望一个真理谓词判定某些语句既真又假，塔斯基 得出结论：除非以某种方式超出该语言的界限，否则我们无法为该语言中的所有语句规定一个真理谓词。换言之，一个语言的真理谓词无法在该语言自身内被定义。

\end{document}